% -*- coding: utf-8 -*-
% !TEX program = xelatex
\documentclass[plain,noanswer,medmath]{../../template/jnuexam}
\usepackage{../../template/kysx}

\begin{document}


\examtitle{2013}{2}


\loadproblems{1}{2013P1}%载入试卷一


\makepart{选择题}{1～8小题，每小题4分，共32分}


\begin{problem}
设$\cos x - 1 = x\sin \alpha(x)$，其中$\left| \alpha(x) \right| < \frac{\pi}{2}$，
则当$x \to 0$时，$\alpha (x)$是\pickout{}
\begin{abcd}
\item 比$x$高阶的无穷小．
\item 比$x$低阶的无穷小．
\item 与$x$同阶但不等价的无穷小．
\item 与$x$等价的无穷小．
\end{abcd}
\end{problem}


\begin{problem}
设函数$y = f(x)$由方程$\cos (xy) + \ln y - x = 1$确定，
则$\lim_{n \to \infty} n\left[f\left(\frac{2}{n} \right) - 1 \right] =$\pickout{}
\begin{abcd}
\item $2$．
\item $1$．
\item $- 1$．
\item $- 2$．
\end{abcd}
\end{problem}


\begin{problem}
设函数$f(x) = \begin{cases}
  \sin x, & 0 \le x < \pi , \\
       2, & \pi \le x \le 2\pi ,
\end{cases}$ $F(x) = \int_0^x f(t) \dt$，则\pickout{}
\begin{abcd}
\item $x = \pi$是函数$F(x)$的跳跃间断点．
\item $x = \pi$是函数$F(x)$的可去间断点．
\item $F(x)$在$x = \pi$处连续但不可导．
\item $F(x)$在$x = \pi$处可导．
\end{abcd}
\end{problem}


\begin{problem}
设函数$f(x) =\begin{cases}
  \frac{1}{(x - 1)^{\alpha - 1}}, & 1 < x < \e, \\
   \frac{1}{x\ln^{\alpha + 1} x}, & x \ge \e,
\end{cases}$若反常积分$\int_1^{+\infty} f(x) \dx$收敛，则\pickout{}
\begin{abcd}
\item $\alpha < - 2$．
\item $\alpha > 2$．
\item $- 2 < \alpha < 0$．
\item $0 < \alpha < 2$．
\end{abcd}
\end{problem}


\begin{problem}
设$z = \frac{y}{x}f(xy)$，其中函数$f$可微，
则$\frac{x}{y}\frac{\pdz}{\pdx} + \frac{\pdz}{\pdy} =$\pickout{}
\begin{abcd}
\item $2yf'(xy)$．
\item $- 2yf'(xy)$．
\item $\frac{2}{x}f(xy)$．
\item $- \frac{2}{x}f(xy)$．
\end{abcd}
\end{problem}


\begin{problem}
设$D_k$是圆域$D = \left\{(x,y)\middle|x^2 + y^2 \le 1 \right\}$在第$k$象限的部分，
记\goodbreak$I_k = \iint_{D_k} (y - x)\dx\dy$（$k = 1,2,3,4$），则\pickout{}
\begin{abcd}
\item $I_1 > 0$．
\item $I_2 > 0$．
\item $I_3 > 0$．
\item $I_4 > 0$．
\end{abcd}
\end{problem}


\useproblem{1}{1}{5}%【同试卷一第一(5)题】


\useproblem{1}{1}{6}%【同试卷一第一(6)题】


\makepart{填空题}{9～14小题，每小题4分，共24分}


\begin{problem}
$\lim_{x \to 0} \left(2 - \frac{\ln(1 + x)}{x} \right)^{\frac{1}{x}} =$ \fillin{}．
\end{problem}


\begin{problem}
设函数$f(x) = \int_{-1}^x \sqrt{1 - \e^t} \dt$，则$y = f(x)$的反函数$x = f^{-1}(y)$
在$y = 0$处的导数$\left. \frac{\dx}{\dy} \right|_{y = 0} =$ \fillin{}．
\end{problem}


\begin{problem}
设封闭曲线$L$的极坐标方程为$r = \cos 3\theta$（$- \frac{\pi}{6} \le \theta \le \frac{\pi}{6}$），
则$L$所围成的平面图形的面积为 \fillin{}．
\end{problem}


\begin{problem}
曲线$\begin{cases}
  x = \arctan t, \\
  y = \ln \sqrt{1 + t^2}
\end{cases}$上对应于$t = 1$的点处的法线方程为 \fillin{}．
\end{problem}


%【类似试卷一第二(10)题】
\begin{problem}
已知$y_1 = \e^{3x} - x\e^{2x}$，$y_2 = \e^x - x\e^{2x}$，$y_3 = - x\e^{2x}$
是某二阶常系数非齐次线性微分方程的$3$个解，
该方程满足条件$y\big|_{x = 0} = 0$，$y'\big|_{x = 0} = 1$的解为$y =$ \fillin{}．
\end{problem}


\useproblem{1}{2}{13}%【同试卷一第二(13)题】


\makepart{解答题}{15～23小题，共94分}


\begin{problem}[points=10]
当$x \to 0$时，$1 - \cos x \cdot \cos 2x \cdot \cos 3x$与$ax^n$为等价无穷小，求$n$与$a$的值．
\end{problem}


\begin{problem}[points=10]
设$D$是由曲线$y = x^{\frac{1}{3}}$，直线$x = a$（$a > 0$）及$x$轴所围成的平面图形，
$V_x$，$V_y$分别是$D$绕$x$轴，$y$轴旋转一周所得旋转体的体积，若$V_y = 10V_x$，求$a$的值．
\end{problem}


\begin{problem}[points=10]
设平面内区域$D$由直线$x = 3y,y = 3x$及$x + y = 8$围成，计算$\iint_D x^2 \dx\dy$．
\end{problem}


\useproblem[points=10]{1}{3}{18}%【同试卷一第三(18)题】


\begin{problem}[points=11]
求曲线$x^3 - xy + y^3 = 1$（$x \ge 0$，$y \ge 0$）上的点到坐标原点的最长距离与最短距离．
\end{problem}


\begin{problem}[points=11]
设函数$f(x) = \ln x + \frac{1}{x}$．
\begin{enumerate}
\item
求$f(x)$的最小值；
\item
设数列$\{x_n\}$满足$\ln x_n + \frac{1}{x_{n + 1}} < 1$，证明$\lim_{n \to \infty} x_n$存在，并求此极限．
\end{enumerate}
\end{problem}


\begin{problem}[points=11]
设曲线$L$的方程为$y = \frac{1}{4}x^2 - \frac{1}{2}\ln x$（$1 \le x \le \e$）．
\begin{enumerate}
\item 求$L$的弧长；
\item 设$D$是由曲线$L$，直线$x = 1,x = \e$及$x$轴所围平面图形，求$D$的形心的横坐标．
\end{enumerate}
\end{problem}


\useproblem[points=11]{1}{3}{20}%【同试卷一第三(20)题】


\useproblem[points=11]{1}{3}{21}%【同试卷一第三(21)题】


\end{document}
